\documentclass[11pt,a4paper]{article}

% ------------------------------------------------------------------ %
%  eitkit — Technical Note (software companion document)
%  Version 0.1.1, July 2026
% ------------------------------------------------------------------ %

\usepackage[american]{babel}
\usepackage[utf8]{inputenc}
\usepackage[T1]{fontenc}
\usepackage{amsmath,amssymb}
\usepackage{booktabs}
\usepackage[margin=2.5cm]{geometry}
\usepackage[numbers,sort&compress]{natbib}
\usepackage{hyperref}
\usepackage{enumitem}
\usepackage{xcolor}

\hypersetup{
  colorlinks=true,
  linkcolor=blue!40!black,
  citecolor=blue!40!black,
  urlcolor=blue!40!black,
}

% ------------------------------------------------------------------ %
\title{\texttt{eitkit} --- Electrical Impedance Tomography\\[4pt]
       Reconstruction Toolkit}
\author{Abd'gafar Tunde Tiamiyu\\[2pt]
        \small Independent Researcher\\[2pt]
        \small\texttt{abdgafartunde@yahoo.com}}
\date{Version 0.1.1 --- July 2026}

\begin{document}
\maketitle

% ------------------------------------------------------------------ %
\section{Overview}

\textbf{Electrical Impedance Tomography (EIT)} reconstructs the interior
conductivity distribution $\sigma(x)$ of a domain $\Omega\subset\mathbb{R}^2$
from boundary voltage measurements. The mathematical problem, posed by
Calder\'on~\cite{Calderon1980}, is severely ill-posed: small perturbations in
boundary data can induce arbitrarily large errors in any
reconstruction~\citep{Mandache2001}.

\texttt{eitkit} is an open-source Python package that implements the complete
\textbf{difference EIT} pipeline:

\begin{center}
\begin{tabular}{@{}ll@{}}
\toprule
Stage          & Implementation \\
\midrule
Mesh generation  & DistMesh2D triangular mesh on the unit disc \\
Forward problem  & P1-FEM with gap electrode model \\
Jacobian         & Adjoint method with sparse LU factorisation reuse \\
Inverse problem  & Tikhonov (L$^2$) + L-curve and TV (L$^1$) via ADMM \\
\bottomrule
\end{tabular}
\end{center}

The toolkit requires only NumPy, SciPy, and Matplotlib. Version~0.1.1 provides
a demo notebook, five synthetic phantom shapes, and 154~unit tests.
The source code, installation instructions, and interactive examples are
available at\\[2pt]
\centerline{\url{https://github.com/abdgafartunde/eit-reconstruction-toolkit}}

% ------------------------------------------------------------------ %
\section{Forward Problem}

\subsection{Governing Equation}

Let $\Omega\subset\mathbb{R}^2$ be a bounded Lipschitz domain. Given a
conductivity $\sigma\in L^\infty(\Omega)$ with $\sigma(x)\ge c>0$, the electric
potential $u\in H^1(\Omega)$ satisfies
\begin{align}
-\nabla\cdot(\sigma\nabla u) &= 0      \quad\text{in }\Omega,
\label{eq:pde}\\[4pt]
\sigma\frac{\partial u}{\partial n} &= j \quad\text{on }\partial\Omega,
\label{eq:bc}
\end{align}
where $j$ is the applied current density with
$\int_{\partial\Omega}j\,ds=0$. A ground condition $u(x_0)=0$ at one
interior node ensures uniqueness.

\subsection{Gap Electrode Model}

Each of the $L$ electrodes is modelled as a single boundary node.
A stimulation pair $(e^+,e^-)$ injects current $+I$ at electrode $e^+$ and
extracts $-I$ at $e^-$. The Neumann load vector $f\in\mathbb{R}^N$ has entries
\[
f_i = \begin{cases}
+I & \text{node $i$ is electrode $e^+$},\\[2pt]
-I & \text{node $i$ is electrode $e^-$},\\[2pt]
0  & \text{otherwise}.
\end{cases}
\]

\subsection{P1 Finite Element Discretisation}

The domain is triangulated with $M$ linear elements on $N$ nodes.
Writing $u=\sum_{i=1}^N u_i\varphi_i$ with P1 hat functions
$\varphi_i$, the weak form yields
\begin{equation}
K(\sigma)\,u = f,
\label{eq:Ksystem}
\end{equation}
where the global stiffness matrix is assembled element-wise:
\[
K = \sum_{e=1}^{M} K^e, \qquad
K^e_{ij} = \frac{\sigma_e}{4A_e}(b_i b_j + c_i c_j), \quad i,j\in\{0,1,2\}.
\]
The gradient coefficients are $b_i=y_{i+1}-y_{i+2}$,
$c_i=x_{i+2}-x_{i+1}$ (indices mod~3), and $A_e$ is the element area.
Assembly is vectorised over elements using NumPy advanced indexing, producing a
\texttt{csr\_array} for efficient sparse arithmetic.

\subsection{Measurement Protocol}

The standard \textbf{adjacent drive / adjacent measurement} protocol is used.
For $L$ electrodes there are $L$ drive steps; step~$k$ sources current at
electrode~$k$ and sinks at $(k+1)\bmod L$. Per drive step, $L-3$ voltage
measurements are taken across adjacent electrode pairs $(k+2,k+3),\dots,(k+L-2,k+L-1)$
(all indices modulo~$L$). The total number of measurements is $P=L(L-3)$; for
$L=16$ this gives $P=208$.

\subsection{Difference EIT and the Jacobian}

Absolute voltages are difficult to calibrate experimentally. Difference EIT
measures voltage changes relative to a reference state $\sigma_0$:
$\delta V = V(\sigma)-V(\sigma_0)$. Linearising around $\sigma_0$ gives the
Born approximation
\begin{equation}
\delta V \approx J(\sigma_0)\,\delta\sigma,
\label{eq:born}
\end{equation}
where $J\in\mathbb{R}^{P\times M}$ is the Jacobian and
$\delta\sigma=\sigma-\sigma_0\in\mathbb{R}^M$ is the conductivity perturbation.

The Jacobian is computed by the \textbf{adjoint method}. For measurement~$i$
associated with drive step~$k$ and electrode pair $(e^+,e^-)$,
\begin{equation}
J_{ie} = -\int_{\Omega_e} \nabla u_k\cdot\nabla\phi_i\,d\Omega,
\label{eq:adjoint}
\end{equation}
where $u_k$ is the forward potential for drive step~$k$ and $\phi_i$ is the
adjoint potential (unit current injected at $e^+$, extracted at $e^-$).
For the adjacent protocol, at most $L$ adjoint solves are needed, and the
factorised stiffness matrix is reused for all solves via
\texttt{scipy.sparse.linalg.splu}.

% ------------------------------------------------------------------ %
\section{Inverse Problem}

\subsection{Tikhonov Regularisation (L$^2$)}

The Tikhonov solution minimises
\begin{equation}
\delta\hat\sigma =
\arg\min_{\delta\sigma}\Bigl\{
\tfrac12\|J\delta\sigma-\delta V\|_2^2
+ \lambda\|\delta\sigma\|_2^2
\Bigr\},
\label{eq:tikhonov}
\end{equation}
with closed-form solution $\delta\hat\sigma=(J^T J+\lambda I)^{-1}J^T\delta V$.
The regularisation parameter $\lambda$ is chosen automatically via the L-curve
method~\citep{Hansen1992}: the optimal $\lambda$ corresponds to the point of
maximum curvature on the log-log plot of residual norm against solution norm.

\subsection{Total Variation (L$^1$ / ADMM)}

TV regularisation promotes piecewise-constant reconstructions:
\begin{equation}
\delta\hat\sigma =
\arg\min_{\delta\sigma}\Bigl\{
\tfrac12\|J\delta\sigma-\delta V\|_2^2
+ \alpha\|D\delta\sigma\|_1
\Bigr\},
\label{eq:tv}
\end{equation}
where $D\in\mathbb{R}^{F\times M}$ is a sparse finite-difference operator on
element-to-element edges. The problem is solved by the Alternating Direction
Method of Multipliers (ADMM)~\citep{Boyd2011} with splitting $z=D\delta\sigma$.
The $\sigma$-update is a quadratic system solved by direct factorisation, the
$z$-update is element-wise soft-thresholding, and the penalty parameter $\rho$
is auto-scaled as $\rho=\operatorname{tr}(J^T J)/\operatorname{tr}(D^T D)$.

% ------------------------------------------------------------------ %
\section{Package Structure}

\begin{verbatim}
eitkit/
  mesh/          DistMesh2D generation, electrode placement
  forward/       P1 FEM assembly, gap model, adjoint Jacobian
  protocol/      Adjacent drive/meas patterns, AWGN noise
  inverse/
    classical/   Tikhonov + L-curve, TV-ADMM
  utils/         Phantom shapes (5), publication-quality plots
\end{verbatim}

The core dependency stack is NumPy $\ge$\,1.24, SciPy $\ge$\,1.10, and
Matplotlib $\ge$\,3.7. Optional extras include PyTorch and JAX for future
learning-based extensions.

% ------------------------------------------------------------------ %
\section{References}

\begin{thebibliography}{10}

\bibitem{Calderon1980}
A.P.~Calder\'on.
\newblock On an inverse boundary value problem.
\newblock In \emph{Seminar on Numerical Analysis and its Applications to
  Continuum Physics}, Rio de Janeiro, 1980.

\bibitem{Mandache2001}
N.~Mandache.
\newblock Exponential instability in an inverse problem for the Schr\"odinger
  equation.
\newblock \emph{Inverse Problems}, 17(5):1435--1444, 2001.

\bibitem{PerssonStrang2004}
P.-O.~Persson and G.~Strang.
\newblock A simple mesh generator in MATLAB.
\newblock \emph{SIAM Review}, 46(2):329--345, 2004.

\bibitem{Hansen1992}
P.C.~Hansen.
\newblock Analysis of discrete ill-posed problems by means of the L-curve.
\newblock \emph{SIAM Review}, 34(4):561--580, 1992.

\bibitem{Boyd2011}
S.~Boyd, N.~Parikh, E.~Chu, B.~Peleato, and J.~Eckstein.
\newblock Distributed optimization and statistical learning via the alternating
  direction method of multipliers.
\newblock \emph{Foundations and Trends in Machine Learning}, 3(1):1--122, 2011.

\bibitem{AdlerLionheart2006}
A.~Adler and W.R.B.~Lionheart.
\newblock Uses and abuses of EIDORS: an extensible software base for EIT.
\newblock \emph{Physiological Measurement}, 27(5):S25--S42, 2006.

\end{thebibliography}

\end{document}
